\begin{frame}{Method \& Setup}
  \centering
  % Progressive residual flow: 1 -> residual -> 2 -> residual -> 4 frames, union below.
\begin{tikzpicture}[
  >=Stealth, node distance=0.7cm,
  stage/.style={draw=PubBlue, fill=PubBlue!85, text=black, rounded corners=3pt,
                minimum width=2.0cm, minimum height=0.95cm, align=center, font=\scriptsize\bfseries},
  res/.style={draw=HidOrange, fill=HidOrange!85, text=black, rounded corners=3pt,
              minimum width=1.55cm, minimum height=0.7cm, align=center, font=\scriptsize\bfseries},
  fin/.style={draw=GoodGreen, fill=GoodGreen!85, text=black, rounded corners=3pt,
              minimum width=3.4cm, minimum height=0.85cm, align=center, font=\scriptsize\bfseries},
  arr/.style={->, thick, gray!70},
]
  \node[stage] (t1) {1 frame\\all faults};
  \node[res, right=of t1] (r1) {residual};
  \node[stage, right=of r1] (t2) {2 frames\\two-phase};
  \node[res, right=of t2] (r2) {residual};
  \node[stage, right=of r2] (t4) {4 frames};
  \node[fin, below=1.05cm of t2] (u) {detected (union)};
  \draw[arr] (t1)--(r1); \draw[arr] (r1)--(t2);
  \draw[arr] (t2)--(r2); \draw[arr] (r2)--(t4);
  \draw[arr] (t1) |- (u);
  \draw[arr] (t2) -- (u);
  \draw[arr] (t4) |- (u);
\end{tikzpicture}

  \vspace{0.6em}
  \begin{columns}[T]
    \begin{column}{0.5\textwidth}
      \footnotesize
      \begin{itemize}
        \item \alert{10} circuits: \eng{ITC'99} + \eng{ISCAS'89}
        \item Collapsed \alert{stuck-at}; one ATPG engine
        \item Non-scan = OpenSTA slack, \alert{bottom 10\%}
      \end{itemize}
    \end{column}
    \begin{column}{0.5\textwidth}
      \footnotesize
      \begin{itemize}
        \item Each stage targets \alert{only the residual}
        \item Backtrack \BTone\,/\,\BThi; \alert{no timeout}
        \item Full-scan baseline: \alert{same engine \& limits}
      \end{itemize}
    \end{column}
  \end{columns}
  \note{The platform is a progressive residual pipeline: one frame on all faults, then two and four frames on the undetected residual, reporting the union over a fixed fault set so coverage is comparable. Setup: ten circuits from ITC'99 and ISCAS'89, collapsed stuck-at, one ATPG engine, non-scan flip-flops are the bottom 10\% by OpenSTA slack, backtrack limits 800 and 5000 with no per-target timeout, and the full-scan baseline uses the same engine and limits for a fair comparison.}
\end{frame}
